% translated-from: en/chapters/04_paper_pipeline.tex@e4a959d4405d3a0ff9d71cc2cef14e592a31610d1bd8bee71993349cd5d0c271
% glossary:        5bca8d57a70f021e20c0af79aa06cd21013495c1a8246887ae889cd8a056f5fd
% model:           claude-opus-4-7

\section{论文生成流水线}\label{sec:paper}

\subsection{流水线概述}\label{sec:paper:overview}

\begin{figure}[t]
  \centering
  \resizebox{\columnwidth}{!}{% F06: paper-generation swim lane — paper-style three-lane workflow.
% Lanes are coloured stage containers; numbered step markers sit to the
% left of each station; the orange refine back-edge is on top.
% style.tikzstyles is loaded by shared/preamble.tex / standalone wrapper.
\begin{tikzpicture}[node distance=1.0cm and 1.6cm]

  % --- six stations laid out in a 3-row grid -------------------------
  \node[station]                                       (i1)  {\figlabel{F6.brainstorm}};
  \node[station, right=2.6cm of i1]                    (i2)  {\figlabel{F6.refine_idea}};

  \node[station emph, below=1.4cm of i1]               (e1)  {\figlabel{F6.bfts}};
  \node[station, right=2.6cm of e1]                    (e2)  {\figlabel{F6.replicate}};

  \node[station, below=1.4cm of e1]                    (w1)  {\figlabel{F6.draft}};
  \node[station emph, right=2.6cm of w1]               (w2)  {\figlabel{F6.review_refine}};

  % --- numbered step markers, perched on the upper-left corner of each
  % station (PaperBench-style badge), so they never overlap any arrow.
  \node[numbered step, above left=0.05cm and -0.2cm of i1.north west] {1};
  \node[numbered step, above left=0.05cm and -0.2cm of i2.north west] {2};
  \node[numbered step, above left=0.05cm and -0.2cm of e1.north west] {3};
  \node[numbered step, above left=0.05cm and -0.2cm of e2.north west] {4};
  \node[numbered step, above left=0.05cm and -0.2cm of w1.north west] {5};
  \node[numbered step, above left=0.05cm and -0.2cm of w2.north west] {6};

  % --- intra-lane forward arrows (thick, paper-style) ----------------
  \draw[big arrow] (i1) -- (i2);
  \draw[big arrow] (e1) -- (e2);
  \draw[big arrow] (w1) -- (w2);

  % --- inter-lane forward arrows: short straight verticals ----------
  \draw[big arrow] (i2.south) -- (e2.north);
  \draw[big arrow] (e2.south) -- (w2.north);

  % --- refine back-edge (orange, on top, label in middle) ----------
  \draw[big arrow, draw=ari-vermilion] (w2.south) to[bend left=35]
       node[edge label, fill=ari-orange!15, font=\sffamily\bfseries\footnotesize, midway]
       {\figlabel{F6.refine}}
       (w1.south);

  % --- lane backgrounds (in background layer) ----------------------
  \begin{scope}[on background layer]
    \node[stage container, fit=(i1)(i2), inner sep=10pt]   (laneI) {};
    \node[stage container, fit=(e1)(e2), inner sep=10pt]   (laneE) {};
    \node[stage container, fit=(w1)(w2), inner sep=10pt]   (laneW) {};
  \end{scope}
  % lane labels OUTSIDE on the right, after the lanes are drawn (top layer)
  \node[stage title, right=4pt of laneI.east, anchor=west] {\figlabel{F6.lane.idea}};
  \node[stage title, right=4pt of laneE.east, anchor=west] {\figlabel{F6.lane.experiment}};
  \node[stage title, right=4pt of laneW.east, anchor=west] {\figlabel{F6.lane.write}};
\end{tikzpicture}
}
  \caption{论文生成泳道:构思、实验、撰写。每条泳道对应一个阶段;
           由 \texttt{review/refine} 回到 \texttt{draft} 的反向
           边是唯一的非单调转移 (\cref{sec:paper:loop})。}
  \label{fig:swim}
\end{figure}

论文生成流水线是 ARI 的第二棵活动树 (\cref{fig:swim})。一旦探
索终止,存活谱系交付给 \term{paper 技能} 生成草稿,再交给
\term{review 技能} 评分。实现是
\code{ari-core/ari/pipeline/orchestrator.py:131} 的
\texttt{run\_pipeline} 函数;同文件第 68 行的
\texttt{build\_scientific\_data} 助手将各节点产物汇集成 paper
技能所需的输入。

一次运行写出三件持久产物:\texttt{draft.tex}(LaTeX 论文本身)、
\texttt{review.json}(按叶的评分细则得分及理由)与
\texttt{rubric.json}(本次运行所用的细则实例,用于动态生成场
景)。

\subsection{评分细则形式}\label{sec:paper:rubric}

我们将\term{评分细则} $\Rubric$ 建模为有限树,叶节点携带元组
$(\rubricitem{i},\rubricweight{i},\judgefn{i})$。每个叶 $i$ 含
类别描述符 $\rubricitem{i}$、非负权重 $\rubricweight{i}$ 满足
$\sum_i \rubricweight{i}=1$,以及\emph{判官函数}
$\judgefn{i}:\paper\to[0,1]$,返回分级得分。论文得分即凸组合
\begin{equation}\label{eq:paperscore}
  \paperscore{\paper} \;=\; \sum_i \rubricweight{i}\,\judgefn{i}(\paper).
\end{equation}
此形式与 PaperBench 的细则树共享,ARI 直接通过 vendor 化的
PaperBench 包(\cref{sec:paper:paperbench})复用之:叶类型
\texttt{TaskNode} 来自 \texttt{paperbench.rubric.tasks},graded
leaf 类型 \texttt{GradedTaskNode} 来自
\texttt{paperbench.judge.graded\_task\_node},加权聚合包络来自
\code{ari-skill-paper-re/src/\_paperbench\_bridge.py:75} 的
\texttt{aggregate\_graded\_tree}。

生产中使用两个判官族:LLM 评分器在 ARI 侧由包装 PaperBench 的
\texttt{SimpleJudge}
(\code{ari-skill-paper-re/src/\_paperbench\_bridge.py:54}) 实
现;确定性评分器用于具备二元或数值校验的项目(如「可复现性」、
「以美元绝对值报告的成本」)。两者均实现 \texttt{Evaluator} 协议
(\code{ari-core/ari/protocols/evaluator.py:19})。

\subsection{review/refine 循环}\label{sec:paper:loop}

改稿循环按反馈 $J(\paper_t)$ 迭代地将 $\paper_t$ 重写为
$\paper_{t+1}$:
\begin{equation}\label{eq:refine}
  \paper_{t+1} \;=\; \refine\!\bigl(\paper_t, J(\paper_t)\bigr).
\end{equation}

当 $\paperscore{\paper_{t+1}}-\paperscore{\paper_t}<\epsilon$
连续两轮成立,或达到流水线级三次迭代上限时,循环终止。每条分支
(收敛 vs.\ 上限)实际触发的频率,是 \cref{chapter:eval} 待冻
结检查点就位后报告的内容;此处不做定量主张。

改稿在每条轴上是有意非单调的:修复清晰度问题可能轻微降低某数值
轴。这正是 \cref{fig:swim} 中反向边的来源;收敛标准作用于聚合
得分,而非按轴。

\subsection{paper-re 的 PaperBench 集成}\label{sec:paper:paperbench}

\term{paper-re 技能}(\texttt{ari-skill-paper-re})是 ARI 的
PaperBench 兼容复现驱动器。集成\emph{不是}重新实现:PaperBench
被 vendor 化于 \texttt{vendor/paperbench},使 ARI 与上游任务与
细则语义保持逐字一致(布局要求记于 PaperBench 的
\texttt{rubric.md \S{}8})。三处集成点值得说明。

\subsubsection*{1. 求解器子类: \texttt{AriPBSolver}}

\code{ari-skill-paper-re/src/\_replicator\_agent.py:269} 子类化
PaperBench 上游的
\texttt{paperbench.solvers.basic.BasicAgentSolver}。chz 不可变
配置规则要求子类亦带装饰器,即便未新增字段。ARI 的覆写最少:

\begin{itemize}
  \item \texttt{\_get\_tools()}:除 submit 外所有工具均被
        \texttt{\_BoundedOutputTool} 包裹,使 bash / python /
        file-read / search 输出在流入下次 API 调用前先大小封顶。
        \texttt{submit} 按名分派、不调用
        \texttt{execute()},因此原样通过。
  \item \texttt{\_run\_agent()}:用
        \texttt{\_bypass\_docker\_sanity\_check} 绕过
        \texttt{paperbench.solvers.utils.sanity\_check\_docker};
        将上游硬编码的 \texttt{/home/\{submission,paper\}} 路径
        改写为 ARI 的工作空间相对布局
        (\texttt{LocalComputer} cwd = \texttt{repro\_sandbox/}、
        论文置 \texttt{paper/}、submission 置
        \texttt{submission/})。提示内容其余从上游
        \texttt{get\_instructions} / \texttt{get\_system\_message}
        逐字而来,使 ARI 与上游复现任务框架(7 天预算、「直至
        复现所有结果不得停止」等)保持一致。
\end{itemize}

\texttt{AriPBSolver} 还保留上游的
\texttt{check\_for\_existing\_run},因此部分完成的运行可幂等地
恢复。

\subsubsection*{2. 细则驱动的产物列表}

vendor PaperBench 不知 ARI 的逐论文细则树,故
\texttt{AriPBSolver.\_run\_agent} 在
\texttt{LocalPBTask.rubric\_expected\_artifacts} 非空时,向用户
指令追加 \texttt{EXPECTED\_ARTIFACTS} 块。该块列出成功的
\texttt{reproduce.sh} 应产出的路径;否则,grader 的叶主张缺乏对
应验证的文件提示。细则树本身由
\code{ari-skill-paper-re/src/\_paperbench\_bridge.py:63}
(\texttt{task\_node\_from\_dict}) 构建。

\subsubsection*{3. 驱动器: \texttt{run\_replicator\_agent}}

\code{ari-skill-paper-re/src/\_replicator\_agent.py:365} 是异步
入口。返回标准 \texttt{build\_reproduce\_sh} 包络
$\{$\texttt{populated},
\texttt{output\_dir}, \texttt{files},
\texttt{expected\_artifacts}, \texttt{max\_runtime\_sec},
\texttt{model}, \texttt{prompt\_sha256}, \texttt{notes},
\texttt{warnings}$\}$。\texttt{output\_dir} 既是智能体的工作
空间,也是 ARI 期望 rollout 后找到 \texttt{reproduce.sh} 的
位置;上游求解器把 \texttt{agent.log} 等簿记写入
\texttt{run\_dir},我们将其指向工作空间本身,使产物与
\texttt{reproduce.sh} 同处。

默认 completer 为 \texttt{OpenAIResponsesTurnCompleterConfig},
模型 \texttt{gpt-5-mini}(可由环境变量
\texttt{ARI\_MODEL\_REPLICATOR} 覆盖);LiteLLM 路由的替代由
\texttt{\_litellm\_completer.py} 助手装配。默认时间上限每次运
行 12 小时;\texttt{sandbox\_kind} 为 \texttt{auto},若提供本
地 Apptainer 镜像则回退至此。

\subsubsection*{4. 评分与聚合}

复现智能体完成后,评分通过
\code{ari-skill-paper-re/src/\_paperbench\_bridge.py} 的
\texttt{judge\_submission} 委派。该函数将 ARI 的
$($\texttt{paper\_md}, \texttt{submission\_dir},
\texttt{reproduce\_log}, \texttt{judge\_model}$)$ 元组适配到上
游 \texttt{SimpleJudge},按 submission 产出
\texttt{GradedTaskNode} 树。运行间聚合由两个助手处理:

\begin{itemize}
  \item \texttt{aggregate\_graded\_tree}(第 75 行):将单个
        \texttt{GradedTaskNode} 通过细则权重 $\rubricweight{i}$
        汇成加权标量,并附按类别细分。
  \item \texttt{average\_graded\_runs}(第 91 行附近):取 $n$
        个 \texttt{GradedTaskNode} 树的逐运行均值,即便智能体
        多次重跑以降噪,也只报告单一复现得分。
\end{itemize}

两阶段设计 —— 叶评分用 vendor 的 \texttt{SimpleJudge},形状与
聚合用 ARI 的适配 —— 意味着升级 PaperBench 上游只是 vendor
swap(无需改 ARI 代码),而 ARI 完全控制如何把叶得分组合成单
一数值。同样的包络在叶为二值时按 $\judgefn{i}\in\{0,1\}$ 流入
\cref{eq:paperscore}。

\subsection{故障模式与护栏}\label{sec:paper:failures}

下述三类故障模式常见到值得命名:

\begin{itemize}
  \item \textbf{LLM-only 支撑}:论文做出谱系中无节点支撑的主
        张。缓解为静态检查:每项数值或宣称的因果主张必须可经
        由 \texttt{\textbackslash code\{...\}} 引用追溯到节点
        产物。
  \item \textbf{引用幻觉}:论文引用并不存在的工作。缓解结构化
        进行:禁止 LLM 杜撰引用键,所有 bib 条目须经 Semantic
        Scholar 校验 (\cref{sec:related})。
  \item \textbf{改稿漂移}:得分上升但论文连贯性下降。缓解为细
        则的按轴下限:任一轴跌破 $0.5$ 即终止改稿。
\end{itemize}

撰稿时的源文件选择由 \code{node\_selection.py:204}
(\texttt{select\_source\_files\_for\_publication}) 完成;它通
过枚举主张可依赖的谱系节点,防止 LLM-only 支撑型故障。
\texttt{load\_selected\_sources}(第 258 行)读取每个
\texttt{(node\_id, rel\_path)} 对的字节,并尊重可选的
\texttt{size\_budget},使论文草稿绝不超出 context。

\subsection{复现器的工具输出边界}\label{sec:paper:bounded-tool}

复现智能体可能在单一 bash 会话中运行数小时;其工具输出(长
stdout、大目录列表)若不处理将主导后续每次提示。ARI 用
\texttt{\_BoundedOutputTool}
(\code{ari-skill-paper-re/src/\_replicator\_agent.py:243}) 加以
界定,该工具包装内部工具的 \texttt{execute()},并经
\texttt{\_truncate\_tool\_output}(第 218 行)按可配置字节上限
后处理输出。这就是 PaperBench 上游的自由形式 transcript 与 ARI
有界、可重放版本的差别;也是多小时 rollout 中提示预算保持可预
测的原因。
